\chapter{The Case for Data Institutions; closing AI's \textit{paradox of openness}}

At present, private firms and AI developers exploit \textit{open AI} discourse and access to the digital commons to deploy publicly available knowledge, tools and data to train models which are then packaged into commercial products. By not returning value to those who generated the data/knowledge, or contributing to the preservation and enrichment of the digital commons, it is depleted (from \textcite{PoliciesDigitalCommons}, in line with \textcite{ostromGoverningCommonsEvolution}). This is AI's \textit{paradox of openness}, made possible by the asymmetrically closed nature of ML development. 

I propose that dedicated DIs for ML datasets close this paradox. Considering a range of technical and policy requirements for ML datasets (from other AI-related literatures), it is possible to give a \textit{fuller} account of the dimensions across which DIs for datasets sustain and enrich the digital commons, as they return financial value to data producers, support pro-social or governance function, and offer additional services. In addition, by introducing mechanisms for democratic accountability along the data value chain, DIs open up novel access points into the inner workings of commercial ML development, incorporate a commitment to a diversity of public values by design, and create value for ML data users in a manner that regulation of models and their outputs \textit{a priori} cannot. This downstream value proposition offsets the increased costs to be potentially faced by data users licensing these datasets. This creates room for manoeuvre in the reassignment of value towards the true costs and benefits around ML datasets. This internalising mechanism moves to close the \textit{paradox of openness}.

\section*{ML's relationship to the digital commons}
\addcontentsline{toc}{section}{ML's relationship to the digital commons}

To unpack how DIs for ML datasets close AI's \textit{paradox of openness}, first we must establish the relationship between ML systems and the digital commons.

\subsection*{Commons as source of representation for ML systems}
\addcontentsline{toc}{subsection}{Commons as source of representation for ML systems}

Continued AI/ML progress is \textit{dependent} on the ongoing health of the digital commons - that is the cumulative digital resources (text, media, code, legal documents, tools) collectively created, shared and managed, with established rules for access and use (\cite{meadwayCREATINGDIGITALCOMMONS}). Most immediately, open-science norms around the distribution of model architectures, together with poor IP protections for software (often considered functional information, over creative expression) result in many models being made freely available over GitHub or Hugging Face as part of a digital commons. However, the dependency is more fundamental still.

The wealth of high-quality information in the digital commons is a principle source of training data for ML systems. ``There is no AI without data" (\cite{groger2021ThereNoAI, 2023WhatWeMean}) and there is no AI without \textit{open} data. Analytically prior to distributional claims as to how the benefits of AI ought to be allocated (to which a commons model has much to say), a robust digital commons is \textit{necessary} to sustain present and future ML development. Recent advances in the field of AI research are a function of increasingly comprehensive and granular representations of our physical and social worlds in machine-readable format - \textit{data}. As increased understanding of empirical scaling laws motivate the use of ever larger amounts of training data (see \textcite{hoffmann2022EmpiricalAnalysisComputeoptimal, kaplan2020ScalingLawsNeural, villalobos2022WillWeRunb}), it is clear that AI systems will continue to draw on the data and knowledge-schema of the commons for representation.\footnote{For technical discussions of the development of an 'AI-ready' semantic web, see \textcite{LinkedDataDesign, WTFKnowledgeGraph}; also from a management perspective see \textcite{criado2019CreatingPublicValue}}

\subsection*{AI threatens the commons}
\addcontentsline{toc}{subsection}{AI threatens the commons}

While AI progress is wholly dependent on a continuously healthy commons, the \textit{paradox of openness} raised by current modes of AI development also threatens its viability. At a technical level, generative AI’s ability to produce synthetic data at scale threatens to pollute the commons (see \textcite{ovadya2019ReducingMaliciousUse}). At an institutional level, curating and maintaining the digital commons as a ‘shared inheritance of mankind’ (\cite{khong2023ARTIFICIALINTELLIGENCECOMMON}) is a sisyphean collective endeavour. Unilateral extraction of value and negation of licensing and IP norms by ML developers destabilises necessary incentives underpinning the commons' delicate self-propagation, degrading quality and paving the way for possible collapse (\cite{meadwayCREATINGDIGITALCOMMONS}). 

Recognising AI governance as existing on a continuum with data governance (\cite{InterwovenRealmsData}), I draw on existing accounts of digital capitalism's extractive logics to provide the contextual dynamics into which AI systems are being introduced (\cite{jameson2019ExtractiveLogicsSystemic}). From surveillance capitalism (\cite{zuboff}), and the ‘data-ification’ of social interactions (\cite{couldry2020CostsConnectionHow, ricaurte2019DataEpistemologiesColoniality}), to the sheer physicality of data infrastructures (\cite{png2022TensionsSouthNorth, 2023PowerEcologyDiplomacy}), redressing the tendencies of digital markets for capturing commons resources has long-formed the backbone of data governance theorising, and appropriate institutional fixes have been roundly considered. AI governance scholarship is not emerging in a vacuum. 

Ricaurte encapsulates the normative ambition of this branch of scholarship in describing the need to address “data-centred economies" which foster extractive models of resource exploitation, human rights violation, cultural exclusion and ecocide, concluding that ``data-extractivism assumes that everything is a data source…life itself is nothing more than a continuous flow of data” \parencite[p. 352]{ricaurte2019DataEpistemologiesColoniality}. AI-enabled business models are set to exacerbate these dynamics. Digital theorists have already identified initial signs of decay within the digital commons as AI developers externalise the costs of generating data while capturing its economic value (\cite{chan2023ReclaimingDigitalCommons}), or as traffic is redirected towards `walled gardens' (\cite{huang2023GenerativeAIDigital}). 

AI governance should leverage lessons learnt from data governance, starting with DIs. Initially theorised to rebalance digital capitalism’s frictionless consumption of the commons, DIs for ML datasets could address these inequality and degradation dynamics built into current modes of training AI systems.

\subsection*{AI indexes the commons}
\addcontentsline{toc}{subsection}{AI indexes the commons}

Cutting against conceptions of AI as commons parasite, the third dimension of this relationship is a vision of AI systems as the digital apotheosis of the commons' current form. AI in this techno-utopian view approaches the summation of human knowledge indexed by algorithms into a ‘blurry Jpeg’ of the internet and query-able by chat boxes (\cite{Chiang_ted}). 

As debatable (and out of scope) as that may be, a positive vision of AI systems as the inheritors of the digital commons needs to first address the \textit{paradox of openness}, such that AI development does not bite the human hand that feeds it. DIs for ML datasets do this.

\section*{How data institutions address AI's \textit{paradox of openness}}
\addcontentsline{toc}{section}{How data institutions address AI's \textit{paradox of openness}}

As strategic objects in ML supply chains, dataset governance gets to the heart the \textit{paradox of openness}. Treating ML datasets as a commons provides a sustainable framework for balancing data use, the rights and interests of data producers, and optimum allocation of public goods across four (overlapping) dimensions. DIs could incorporate a dataset-as-a-service approach alongside community accountability and input to deliver higher quality \textit{datasets}, as well as higher quality data \textit{usage} around risks, liability, licensing and consent. At a governance level, the distributed leverage of an ecosystem of DIs would empower and incentivise creators, promoting greater \textit{sustainability}, alongside \textit{equity and inclusion}. In this way, DIs as commons-based data governance institutions incorporate the supply of public goods into the ML value chain, offset by creating new dimensions of value for downstream users.

\subsection*{Data quality}
\addcontentsline{toc}{subsection}{Data quality}

Beyond DIs as a  resource-to-infrastructure (data-to-dataset) layer upon which to build the goods and services of an AI economy, I propose that their continuous and iterative provision of \textit{quality} datasets is infrastructural to a healthy digital society and body politic. Delivering quality is demanding work - mostly unnoticed, until it isn't. Without adequate opportunities for scrutiny and accountability, ML data markets under-allocate quality.

There exists an extensive literature as to what quality in data means, for whom, towards which application, and how its dimensions can be measured and managed (\cite{gudivada2017DataQualityConsiderationsa, munappy2022DataManagementProduction, sambasivan2021EveryoneWantsModel, kreutzer2022QualityGlanceAudit}), I operationalise quality datasets as a function of robust data work, being both \textit{safety enhancing} (debiased, de-risked, well-sampled, coherently licensed, consensual, audited, etc) and \textit{ML-ready} (methodically annotated, balanced across classes, low bias, feature-rich, well-structured, etc) - a composite definition drawn from \textcite{paullada2021DataItsDis} and \textcite{gudivada2017DataQualityConsiderationsa}. Quality is value-agnostic (as is the DI institutional template) since the `appropriate values' are use case dependent, and quality provides value both to upstream communities, and to downstream users (developers as well as future data subjects).

The expected quality gains derived from participation in open data projects is an area of active empirical research (\cite{criado2019CreatingPublicValue, 2023ExploringValueDecentralised, 2021EconomicImpactTrust, }). Open data \textit{in principle} allows for a representative range of social values to be incorporated and for community feedback and iteration, but also risks data pollution, value divergence, and community neglect (\cite{2022MeasuringImpactData}). DIs navigate this dilemma by providing institutional templates to democratically establish context-appropriate practices for data stewardship. Specifically, curating quality datasets is a difficult, expensive and thankless process - even more so the next generation of rigorously labelled, metadata rich, ML-ready datasets (\cite{oala2023DMLRDatacentricMachine}). The true social cost of data work is simultaneously expensive and undervalued. ML developers are incentivised to free-ride on quality to capture market share, such that quality is under-allocated by the market. This results in incomplete, noisy datasets characterised by representational harms, biases, poor encoding, and imbalanced distributions (\cite{gudivada2017DataQualityConsiderationsa}). 

Ultimately, companies are not designed to arbitrate public good in politically-contested spaces at scale. Companies are a legal fiction to avoid liability and maximise shareholder value. The qualities society requires of its digital infrastructure do not necessarily align with those prioritised by private firms, and analyses of data supply chains find compromised quality, obfuscated traceability and routine accountability-avoidance (\cite{longpre2023DataProvenanceInitiative}). As  artefacts functionally-dedicated to supplying the governance and technical-infrastructure for developing quality datasets, DIs offer the possibility to embed community-principles and selected public values into the design of the datasets they steward. ML data management practitioners emphasise that identifying and funding development of safety-enhancing and ML-ready datasets - among other human-design principles (see \textcite{bai2022TrainingHelpfulHarmless}) - requires ongoing consultation with diverse communities, impacted stakeholders, and experts across disciplines (\cite{aragon2022HumanCenteredDataScience}), as well as transparency and ongoing ethical consideration (\cite{jo2020LessonsArchivesStrategies}). DIs provide ideal fora to facilitate the necessary engagement. In this sense, DIs double as quality-guarantors that a dataset is high-quality, human-generated, and suitable for its stated purposes.

\subsection*{Risk, liability, licensing, consent}
\addcontentsline{toc}{subsection}{Risk, liability, licensing, consent}

Currently ML developers freely extract value from openly available commons data without adequately pricing-in associated harms, taking liability or offering recourse. This is particularly true of web scraping publicly available informational sources. The result is the externalisation of the risk and social costs associated with ML development alongside the private capture of commons value. DIs close this dimension of the \textit{paradox of openness} by incentivising better data usage. 

\subsubsection*{Compliance, liability \& delegated licensing}
\addcontentsline{toc}{subsubsection}{Compliance, liability \& delegated licensing}

Building on the emphasise on safety and internalisation of data’s true costs in the definition of quality above, DIs assuming responsibility for collecting the consent of impacted communities and individuals, implementing necessary privacy protections, and able to guarantee regulatory compliance, would represent a gain for ML data users (\cite{chan2023ReclaimingDigitalCommons}).

A network of representative stewarding institutions empowered to manage and license data to commercial ML developers in a manner consistent with the consent provided, would be further well-positioned to provide clarity over licensing options for divers ML use cases. The economics of licensing data for training LLMs and other sophisticated data-intensive modes remains intractable and beyond scope, but DIs provide an inclusive basis from which to begin an informed public discussion. Further, by collectively offering a gradation of tailored policy options differentiated by risk tolerance or other domain-specific sets of preferences, the `market for DIs' allows individuals to select that which is most appropriate to their individual circumstances (\cite{2021WhatAreBottomup}). In aggregate, this sustainably incentivises greater data sharing and alleviates a source of friction in AI development.

A related constraint when opening datasets for ML training is the need to include harmful content for model learning (to be able to classify such content as harmful in newly sampled data). It remains problematic to make known harmful content open-access; DIs offer a partial solution. As a function of their proximity to the data-generating community, DIs could manage a separate category of data, evaluating toxicity (\cite{gehman2020RealToxicityPromptsEvaluatingNeural}) and truthfulness (\cite{lin2022TruthfulQAMeasuringHow}) as it pertains to the community impacted. While not a complete solution, there’s no free lunch in statistics or policy. Collective data governance can maximally balances accountability and control over harmful content against greater sharing, with a legitimacy that current modes of ML dataset curation lack.

\subsubsection*{Supporting opt-out mechanisms}
\addcontentsline{toc}{subsubsection}{Supporting opt-out mechanisms}

DIs could better support opt-out mechanisms for privacy, benefitting both data subjects and model developers’ business models, risk and approval ratings. 

At present, data creators cannot prevent AI developers from using their data. Firms are by design minimally accountable under the law, opt-out mechanisms are lacking, and training datasets are closed, such that even if an individual were to threaten to withhold their data from a model developer, they lack knowledge (as to who is using their data) and bargaining power (\cite{chan2023ReclaimingDigitalCommons}).

In contrast, direct community accountability incentivises DIs to facilitate individuals’ data-forgetting rights by promoting technical traceability, provenance documentation and opt-out mechanisms in the datasets they steward (see \textcite{chan2023ReclaimingDigitalCommons} for technical proposals). Further, as contracting organisations DIs can \textit{license} the data they steward. When licensing data, DIs could require transparent processes from model developers for removing the influence of individual data points (see \textcite{nguyen2022SurveyMachineUnlearning}), or provide a coordinating hub for the erasure of an individual’s data across all models where it is used. DIs thus empower producers to grant permission for their data to be used in one context but not in another, and to determine the after-life of the content to which they contribute. 

\subsection*{Sustainability}
\addcontentsline{toc}{subsection}{Sustainability}

Data-centric AI scholarship calls for data to be reconceptualised as an ever-evolving component in an AI system, to be iteratively augmented and developed across the system’s lifecycle (\cite{jarrahi2023PrinciplesDataCentricAI}). This view emphasises that ML developers require dependable supply chains providing up-to-date high-quality data at scale, continuously. Meeting this demand for data sustainability requires managing the delicate balance between openness and safeguarding the rights and interests of contributors. This is not-possible with blunt binary-choice \textit{open AI} policy tooling.

DIs offer gradations of openness and return value (financial or other) to data producers. By realigning dataset development incentives closer to the needs of data-generating communities, DIs create more than just a legal justification and infrastructure for high-volume data sharing, rather they represent organisations worthy of societal trust (\cite{mcdonaldReclaimingDataTrusts}). It is trust that sustains the digital commons, and the AI systems it supports. DIs thus move towards closing this dimension of AI's \textit{paradox of openness} by offering a governance model for sustainable data sharing infrastructure.

\subsection*{Equity \& inclusion}
\addcontentsline{toc}{subsection}{Equity \& inclusion}

The final harmful dimension associated with AI's \textit{paradox of openness} that I propose DIs could be leveraged to address, is the reality that dominant global discourses around openness and inclusion achieve contrary aims. This happens through satisficing and non-representative inclusion. Here I draw on divers critical literatures (such as post/decolonial computing, digital divide research, data justice, indigenous data sovereignty, and cultural pluralism in AI - among others) to propose that a decentralised network of bottom-up DIs offer constructive solutions to the open questions and critiques raised by these research programmes. 

\subsubsection*{The paradox of participation within \textit{inclusive AI}}
\addcontentsline{toc}{subsubsection}{The paradox of participation within \textit{inclusive AI}}

Narrow \textit{open AI} proposals focus on making AI technologies and resources available to a broader audience, alongside facilitating collaboration and the free exchange of knowledge and data. This preoccupation with access can frustrate democratic outcomes, as unstructured access to AI systems prevents societies from restricting those uses they deem undesirable (\cite{chan2023ReclaimingDigitalCommons}). Further, in practice developers instrument \textit{open AI} discourses to maintain open access to the commons while expropriating and repackaging communal value (\cite{graywidder2023OpenBusinessBiga}). It is not enough to simply make AI products freely available (\cite{PoliciesDigitalCommons}). Parallel to this \textit{open AI} dynamic, paying lip service to \textit{inclusive AI} discourses (of which there are many variants) can similarly result in exclusion, if not accompanied by intentional efforts to include underrepresented groups (\cite{png2022TensionsSouthNorth}).

Openness (\textit{open AI}) and numerical inclusion (\textit{inclusive AI}) alone will not dismantle existing power imbalances in AI governance. Instead policy makers must by attentive to the structures of power and systems of values that are imprinted onto technology’s design and development, and be intentional as to whose representation gets included within that process. Tools are not value-agnostic, and the view from nowhere is nowhere to be found. Democratising AI involves dismantling and reassembling mechanisms (technical, institutional, or otherwise) within AI development practices that perpetuate inequalities, and DIs offer a template for institutionalising these considerations.

Divers critical literatures highlight the \textit{paradox of participation} at the heart of global AI governance (\cite{chatham_house}). Critics variously point to the reality of its domination by Global North institutions (\cite{jobin2019GlobalLandscapeAI, png2022TensionsSouthNorth}), its neglect of local knowledge, cultural pluralism and the demands of global fairness (\cite{singh2021MappingAIGlobal}), its failure to develop adequate protocols to redress capture by vested interests (\cite{ulnicane2021GoodGovernanceResponse}), and to processes of ‘predatory inclusion’ at the all levels of global governance - intended for procedural, numeric, and promotional purposes (\cite{bliss2008ParticipationInternationalDevelopment}). Hierarchy-preserving multi-stakeholder governance initiatives such as GPAI, UN Tech Envoy Office, UNESCO, WEF, IEEE etc, are seen to be proffered in lieu of structural change, driving the \textit{paradox} wherein inclusion can exist while structural harms persist (\cite{png2022TensionsSouthNorth}). This ‘violence of care’ is intended to satisfice disenfranchised stakeholders and minimise friction in advancing the status quo-preserving trajectory (see \textcite{cleaver1999ParadoxesParticipationQuestioning} for initial theorising). The result of this governance dysfunction are classes of learning algorithms trained on London’s traffic flows or public health datasets, that with some contextual fine-tuning are ready to bring yet more benefit to Washington or Sydney-based stakeholders, yet which remain uninformative for the streets or hospitals of Kampala or Jakarta where material conditions diverge abruptly  (\cite{smith2018ArtificialIntelligenceHuman}). The cycle continues and \textit{inclusive AI}-style rhetoric perpetuates the further bifurcation of an ‘AI Divide’ (\cite{lindroth2017GlobalPoliticsIts}). 

DIs offer a possibility for disenfranchised stakeholders to resist this dynamic of AI systems compounding long-standing geopolitical inequalities (\cite{couldry2020CostsConnectionHow}). As a function of their grassroots representation of data-generating communities and the leverage they offer for withholding data, DIs give credence to calls for localised AI regulation and development operationalised to their immediate cultural and infrastructural contexts (\cite{kak2020GlobalSouthEverywhere}). 

\subsubsection*{Digital Sovereignty}
\addcontentsline{toc}{subsubsection}{Digital Sovereignty}


In contrast to \textit{inclusive AI}-style optimism around `leapfrogging' and `digital transformation' (\cite{singhDigitalIndustrialisationDeveloping, borokini2021ExploringthefutureofdatagovernaceinAfricaDataStewardship}), Global South digital development concerns frequently foreground digital sovereignty. DIs in this context serve as possible vectors of resistance to perceived data colonialism by institutionalising self-representation on communities' own terms.

Heterogeneity defines post-colonial scholarship. Nevertheless, Global South concerns around AI coalesce around concerns about data capitalism/colonialism (see \textcite{couldry2020CostsConnectionHow, ricaurte2019DataEpistemologiesColoniality, graham2014CriticalPerspectivePotential}), about dependency dynamics (deriving from standard organisations' regulatory monopolies, unequal Public-Private Partnerships, or global IP regimes) (\cite{png2022TensionsSouthNorth}), and concerns about harms associated with the labour and material supply chains of AI technologies (\cite{2023PowerEcologyDiplomacy}). These concerns have led to documented core-periphery patterns of under-representation and data availability in ML datasets (\cite{graham2014CriticalPerspectivePotential}), leading to distributional shifts, generalisation errors (even conceptual invalidity), and contextual incompatibilities in imported AI systems (\cite{smith2018ArtificialIntelligenceHuman}).

As grassroots democratic institutions capable of withholding the means of AI production,  DIs allow communities to dictate their own locally-appropriate patterns of engagement with globalised ML value chains. In the context of the Global South, this empowerment could encourage greater data sharing amongst long-suspicious populations (see \textcite{olorunju2023AfricanDataTrusts} for a theoretical proposal, and \textcite{baarbe2019ProposedAgriculturalData} for case study application to the field of agriculture). To the data management literature's concerns, DIs support intentionally collected datasets rooted in their original contexts, distributed in ways that respect the privacy, property and rights of data creators, and constructed in conversation with impacted stakeholders or domain experts (\cite{paullada2021DataItsDis}). Taken as a whole, DIs and the digital sovereignty they support chart a much-needed inclusive path forward for AI governance writ large.

\subsubsection*{“Nothing about us, nothing without us”}
\addcontentsline{toc}{subsubsection}{“Nothing about us, nothing without us”}

DIs support the participatory action principle that no policy should be developed without the direct participation of affected communities (see \textcite{mcdowell2023RethinkingArtificialIntelligence}). As representation institutions, DIs centre the knowledge of those communities most exposed to risks, and provide accountability fora for impacted stakeholders to scrutinise and surface hitherto unrecorded design decisions in ML datasets (\cite{frey2020ArtificialIntelligenceInclusion}). In this way, DIs institutionally embody core `participatory AI' principles that seek to involve a wider range of stakeholders than just technology developers and product risk assessors (see \textcite{ortega2020ManagingTransitionMultistakeholder} for an agenda). DIs provide a template for institutionalising guidance as to what ML tasks (social or technical) are valid to a given community, and the infrastructure to enable creation of datasets to support those tasks. 

Further, by supporting grassroots democratic participation in the ML development process, DIs could assemble communities structured by shared interests and provide a mechanism for their preferences to be aggregated into leverageable representation across higher-level advocacy platforms. Rather than subsuming stakeholders into existing top-down global governance structures without making the necessary adjustments to enact meaningful change (\textit{inclusive AI}), DIs build from the bottom-up, accounting for heterogeneous requirements along the way. 

In this way DIs could serve as an institutional response to \textit{reality-centric AI} critiques put forward by \textcite{RealityCentricAI}. This view advocates for datasets that not only advance ML as a field, but contribute positively to society. Advocates highlight the liberatory potential of intentional dataset design to make visible diverse patterns of injustice in the world, such that they may be addressed (see \textcite{DB4L} for an example of reality-centric dataset development). DIs for ML datasets institutionalise this line of reasoning. 

